%==============================================================================
% Chapitre 16 - Le canon
%==============================================================================

\section{Introduction}

Le canon est l'un des composants les plus importants d'une arme à feu.

Son rôle principal consiste à guider le projectile pendant sa phase
d'accélération tout en contenant les gaz de combustion générés par la charge
propulsive.

Les caractéristiques du canon influencent directement :

\begin{itemize}
    \item la vitesse initiale ;
    \item la précision ;
    \item la stabilité du projectile ;
    \item la durée de vie de l'arme ;
    \item les performances balistiques globales.
\end{itemize}

\begin{aretenir}

Le canon est l'interface principale entre l'énergie produite par la munition
et le projectile.

\end{aretenir}

\section{Fonctions du canon}

Le canon assure plusieurs fonctions simultanément :

\begin{itemize}
    \item résister à de fortes pressions ;
    \item guider le projectile ;
    \item assurer l'étanchéité des gaz ;
    \item communiquer une rotation au projectile dans le cas des canons rayés.
\end{itemize}

Ces fonctions doivent être assurées durant un intervalle de temps extrêmement
court, généralement de l'ordre de quelques millisecondes.

\section{Constitution générale}

Un canon moderne est généralement composé :

\begin{itemize}
    \item d'une chambre ;
    \item d'une âme ;
    \item d'une bouche ;
    \item éventuellement d'organes ou accessoires associés.
\end{itemize}

La chambre reçoit la cartouche.

L'âme constitue le conduit parcouru par le projectile.

La bouche correspond à l'extrémité de sortie.

\section{L'âme du canon}

L'âme est la surface intérieure du canon.

Elle guide le projectile pendant son déplacement.

Sa géométrie influence directement :

\begin{itemize}
    \item les frottements ;
    \item l'étanchéité ;
    \item la précision ;
    \item l'usure.
\end{itemize}

Deux grandes familles existent :

\begin{itemize}
    \item les canons lisses ;
    \item les canons rayés.
\end{itemize}

\section{Canons lisses}

Dans un canon lisse, l'âme ne comporte pas de rayures.

Le projectile ne reçoit donc pas de rotation significative due au canon.

Cette solution est utilisée notamment pour :

\begin{itemize}
    \item certains fusils de chasse ;
    \item certains canons de chars modernes ;
    \item certaines armes spécialisées.
\end{itemize}

\section{Canons rayés}

La majorité des armes militaires modernes utilisent un canon rayé.

Des rainures hélicoïdales sont usinées à l'intérieur du canon.

Ces rayures imposent une rotation au projectile pendant son déplacement.

Cette rotation améliore considérablement la stabilité du vol.

\begin{aretenir}

Les rayures permettent de stabiliser le projectile par effet gyroscopique.

\end{aretenir}

\section{Champs et rayures}

Dans un canon rayé, on distingue :

\begin{itemize}
    \item les champs ;
    \item les rayures.
\end{itemize}

Les champs correspondent aux parties les plus proches du diamètre nominal.

Les rayures sont les zones creusées.

La différence entre les deux définit la profondeur des rayures.

\section{Pas de rayure}

Le pas de rayure décrit la vitesse de rotation imposée au projectile.

Il est généralement exprimé sous la forme :

\begin{quote}
1 tour en X pouces
\end{quote}

ou :

\begin{quote}
1 tour en X millimètres
\end{quote}

\begin{exemple}

Un pas de :

\begin{equation}
1:12''
\end{equation}

signifie que le projectile effectue un tour complet après avoir parcouru
12 pouces dans le canon.

\end{exemple}

Le choix du pas dépend principalement :

\begin{itemize}
    \item de la longueur du projectile ;
    \item de sa masse ;
    \item de sa vitesse ;
    \item de son aérodynamique.
\end{itemize}

\section{Longueur du canon}

La longueur du canon influence directement l'accélération du projectile.

Dans de nombreux cas :

\begin{itemize}
    \item un canon plus long permet une vitesse initiale plus élevée ;
    \item un canon plus court réduit l'encombrement.
\end{itemize}

Cette relation possède toutefois des limites.

Lorsque la pression des gaz devient insuffisante, augmenter davantage la
longueur du canon n'apporte plus de gain significatif.

\section{Matériaux}

Les canons modernes sont généralement fabriqués à partir d'aciers à haute
résistance.

Ils doivent supporter :

\begin{itemize}
    \item des pressions élevées ;
    \item des températures importantes ;
    \item des sollicitations répétées.
\end{itemize}

Certains canons reçoivent des traitements particuliers :

\begin{itemize}
    \item chromage ;
    \item nitruration ;
    \item revêtements spécialisés.
\end{itemize}

\section{Échauffement}

Lors du tir, une partie importante de l'énergie est transformée en chaleur.

Le canon s'échauffe progressivement.

Cette élévation de température peut provoquer :

\begin{itemize}
    \item une modification des dimensions ;
    \item une dégradation de la précision ;
    \item une augmentation de l'usure.
\end{itemize}

Ce phénomène est particulièrement important pour les armes automatiques.

\section{Usure}

L'usure du canon résulte principalement :

\begin{itemize}
    \item des hautes températures ;
    \item des fortes pressions ;
    \item du frottement des projectiles ;
    \item de l'érosion chimique.
\end{itemize}

Les zones les plus sollicitées sont généralement :

\begin{itemize}
    \item la chambre ;
    \item la prise de rayures ;
    \item la bouche.
\end{itemize}

\section{Influence sur la précision}

Le canon joue un rôle majeur dans la précision d'une arme.

Parmi les facteurs importants figurent :

\begin{itemize}
    \item sa rigidité ;
    \item sa qualité de fabrication ;
    \item son état d'usure ;
    \item sa température ;
    \item la régularité des rayures.
\end{itemize}

\section{Influence sur la balistique}

Les caractéristiques du canon influencent directement :

\begin{itemize}
    \item la vitesse initiale ;
    \item la rotation du projectile ;
    \item la stabilité du vol ;
    \item la dispersion ;
    \item les performances terminales.
\end{itemize}

De nombreux phénomènes étudiés dans les parties suivantes trouvent leur
origine dans les caractéristiques du canon.

\begin{simulation}

Dans les simulations simplifiées, le canon est souvent représenté par sa
longueur et la vitesse initiale obtenue.

Les modèles avancés peuvent intégrer :

\begin{itemize}
    \item le pas de rayure ;
    \item l'usure ;
    \item l'échauffement ;
    \item les vibrations.
\end{itemize}

\end{simulation}

\section{Vers les chapitres suivants}

Les chapitres consacrés aux projectiles réutiliseront plusieurs notions
introduites ici :

\begin{itemize}
    \item rotation ;
    \item stabilité gyroscopique ;
    \item pas de rayure ;
    \item interaction projectile-canon.
\end{itemize}

\section{À retenir}

\begin{aretenir}

\begin{itemize}
    \item Le canon guide et accélère le projectile.
    \item Les rayures communiquent une rotation stabilisatrice.
    \item Le pas de rayure influence directement la stabilité.
    \item La longueur du canon influence la vitesse initiale.
    \item L'échauffement et l'usure peuvent affecter les performances.
\end{itemize}

\end{aretenir}
